\documentclass[12pt]{article}


%=============================================
% --- Packages ---
%=============================================
% --- Page size
\usepackage[top=1in, bottom=1in, left=1.5in, right=1.5in]{geometry}

%=============================================
% --- Extra color package that enables table coloring and other functionalities
\usepackage[dvipsnames, table, xcdraw]{xcolor}  

%=============================================
% --- Font and Typography ---
\usepackage[utf8]{inputenc}        
\usepackage{sans} 
\usepackage{pifont} 

%=============================================
% --- Graphics and Hyperlinks ---
\usepackage{graphicx}              
\usepackage{caption,subcaption}    
\usepackage{hyperref}              
\hypersetup{                       
	colorlinks=true,               
	linkcolor=blue,                
	filecolor=magenta,             
	urlcolor=cyan,                 
	citecolor=ForestGreen          
}
\graphicspath{{Figures/PDFs}}

%=============================================
% --- Math packages ---
\usepackage{amsmath}               
\usepackage{amssymb}               
\usepackage{amsfonts}              
\usepackage{amsthm}                
\usepackage{mathtools}
\usepackage{array}
\usepackage{multirow}
\usepackage{booktabs}


%%=============================================
% --- List and Box Customization ---
\usepackage{enumitem}  % Customizable lists
%\usepackage{tcolorbox} % Colored boxes
\usepackage[breakable,skins]{tcolorbox}

\theoremstyle{style-one}
\newtheorem{definition}{Definition}[section]  	% Definition environment
\newtheorem{proposition}{Proposition}[section]  % Proposition environment
\newtheorem{theorem}{Theorem}[section]  		% Theorem environment
\newtheorem{example}{Example}[section]          % Example environment
\newtheorem{exercise}{Exercise}[section]        % Exercise environment
\newtheorem{lemma}{Lemma}[section]              % Lemma environment
\newtheorem{problem}{Problem}[section]          % Problem environment


% --- Unnumbered Theorems --- 
\newtheorem*{definition*}{Definition}  % Unnumbered definition environment
\newtheorem*{corollary}{Corollary}     % Unnumbered corollary environment
\newtheorem*{remark}{Remark}

\newtheorem*{problem*}{Problem}  % Unnumbered problem environment


%=============================================

%=============================================
% --- New commands ---
\newcommand{\uP}{u}
\newcommand{\sP}{\mathbf{s}}
\newcommand{\eP}{\mathbf{e}}
\newcommand{\lP}{\lambda}
\newcommand{\mP}{\boldsymbol{\mu}}

\newcommand{\stP}{\mathbf{\tilde{s}}}
\newcommand{\etP}{\mathbf{\tilde{e}}}


\newcommand{\dx}{d\mathbf{x}}
\newcommand{\CG}{\text{CG}}
\newcommand{\VCG}{[\CG]^d}
\newcommand{\DG}{\text{DG}}
\newcommand{\VDG}{[\DG]^d}
\newcommand{\BDM}{\text{BDM}}
\newcommand{\RT}{\text{RT}}
\newcommand{\Hdiv}{\text{Hdiv}}




%=============================================
			  % MAIN DOCUMENT 
%=============================================
\title{Numerical approximation of an optimization problem that appears in  data-driven Computational Mechanics}

\date{\the \year}

\begin{document}

\maketitle

\begin{abstract}
	In this work, we investigate an optimization problem that appears when dealing with the paradigm of Data-Driven Computational Mechanics. In the context of the diffusion-reaction problem, such an optimization problem seeks for the set of continuous primal fields (i.e., a potential field, its gradient and the flux) that satisfies the physical conservation law and the geometrical compatibility at fixed discrete fields (i.e., reconstructed gradient and flux that locally look like the data sets). Moreover, we consider a reaction term in the physical conservation law that automatically couples all the optimality conditions and thus, adding an interesting twist to the problem. First, we establish the well-posedness in the continuous setting and second, we propose a finite-element discretization that approximates the underlying continuous problem while preserving the saddle-point structure and improving the stability properties when compared to plain discretizations. Finally and to substantiate our claims, we present a set of examples that illustrates the advantages of the proposed approach.
\end{abstract}

\textbf{Keywords:} data-driven computational mechanics, optimization problem, diffusion-reaction problem, stabilized finite elements. 

%=============================================
% --- Introduction --- 
\section{Introduction}

\textit{Data-Driven computational mechanics} (DDCM) represents the latest paradigm in continuum mechanics, which can, of course, be employed to a large range of applications. For instance, in the context of the diffusion-reaction problem, the gradient of a concentration scalar field, corresponding to a conserved quantity, and the flux vector field over a physical domain are then to be found as the closest to some field extension of a given finite set of empirically determined gradient-flux pairs. Such a problem can be stated as an optimization problem in which discrete and continuous fields are to be simultaneously sought such that the physical conservation law and the compability condition are satisfied in some weak sense. In this context, a more basic question on the existence of the continuous fields provided that the discrete fields are fixed and given \textit{a priori} naturally arises. This ``optimization subproblem'' is actually a continuous optimization problem that can, in principle, be treated computationally and analyzed using standard tools from the classical theory of linear \textit{partial differential equations} (PDEs). Therefore, answering that question represents the main target of the present contribution. To place the proposed ideas into an appropriate setting, we provide below some of the most important contributions to present date, highlight the difference of the present work with respect to existing literature and state the novelty from the methodological point of view.

DDCM is a very active field of research whose first contributions can be traced back to less than a decade ago and has mostly been developed within solid/structural mechanics. For instance, this principle is applied to the static analysis of multi-bar structures \cite{Kirchdoerfer2016}. Based on similar ideas, a data-driven magnetostatic solver is proposed in \cite{Galetzka2020}. In the following, this initial approach is modified to improve robustness due to the presence of outliers in the constitutive data \cite{Kirchdoerfer2017}. Finally, this is extended to the dynamic analysis of structures that exhibit configuration-independent mass matrices \cite{Kirchdoerfer2018}. For the sake of completeness, we must mention that there are some efforts addressing the existence of solutions for data-driven elasticity \cite{Conti2018,Conti2020}. There is also a hybrid approach that relies on a smooth constitutive manifold reconstructed from the data at a off-line pre-step. This principle is applied to time-independent problems \cite{Gebhardt:2020a}. Then, the approach is also modified to handle time-dependent problems \cite{Gebhardt:2020b}. Lastly, this is extended to deal with \textit{mathematical programs with complementarity constraints} (MPCC) such as those arising in contact mechanics  \cite{Gebhardt2024a}. Recently in  \cite{Gebhardt2024b}, the structural properties of the DDCM problem discretized by means of the \textit{finite element method} (FEM) is investigated and its global solvability is established. Remarkably, there are neither antecedents of DDCM in the context of the diffusion-reaction problem nor direct connections to PDE-regularized methods to analyze empirical data \cite{Ponti2021} or to inhomogeneities identification from boundary measurements \cite{Ammari2004}.

All the works introduced above (that deal with the numerical solution of the discrete-continuous optimization problem in DDCM) require along the solving process to compute the solution of a type of continuous optimization subproblem that is a particular case of the optimization problem addressed in the present work. This is due to that those works adopt an \textit{alternating direction method} (ADM). This work can be considered then as a first step towards the application of DDCM in the context of the diffusion-reaction problem, which to the best of our knowledge has not been realized so far. In previous works, where once the assignation of the data pairs has been done, the optimization subproblem provides two sets of optimality conditions that are totally decoupled and can be solved independently. In the present work, we consider a reaction term present in the physical conservation law that automatically couples all the optimality conditions and thus, the well-posedness of this rather more general optimization problem ought to be established right from the beginning. Moreover, our analysis considers two- and three spatial dimensions in contrast to those more practical contributions introduced before that focus on problems defined on a closed interval of the real line \cite{Gebhardt:2020a, Gebhardt:2020b, Gebhardt2024a, Gebhardt2024b}. In addition, we also introduce a new finite element formulation, establish its well-posedness for the chosen finite element spaces and show its consistency with respect to the continuous formulations posed in the setting of infinite-dimensional Hilbert spaces. The desired stability properties are achieved by directly approximating the saddle-point problem, approach that exhibits some connection to other well-established approaches \cite{Masud2002, Bochev2006, Badia2010, Burman2023}. Moreover, we also consider the case in which the data sets are generated from a Darcy-Forcheimer flow type \cite{Girault2008}.

The remaining of this article is organized as follows. In Section 2, we describe the optimization problem in the reaction-diffusion context. Relying on the theory for quadratic optimization problems subject to linear constraints, we prove that this optimization problem in its continuous setting is indeed solvable with a unique solution. Moreover, we also introduce its spatial approximation by adopting a finite element approach. Then, we prove that the proposed approximating problem is also well posed and consistent. In Section 3, we introduce the numerical implementation of the spatially discretized problem, which is realized by its coding within the Firedrake finite element platform. The resulting code is then briefly explained and made publicly available. This development is followed as well by a series of examples, of increasing difficulty, that illustrate the remarkable properties of the setting proposed. Finally, in Section 4, we conclude with remarks, potential limitations and possible future work.  


%=============================================
% --- Theory ---
\section{The model problem}

We consider a reaction-diffusion problem governed by the conservation law

\begin{equation}
\nabla \cdot \sP + \zeta\,\uP = q
\label{eq:conservation}
\end{equation}

posed on a regular domain $\Omega \subset \mathbb{R}^{d_{\Omega}}$. In equation (\ref{eq:conservation}) the symbol $\nabla$ represents the vector of partial derivatives, so that $\nabla \cdot \sP$ is the divergence of the vector field $\sP$, which represents the flux. Also, $\uP$ is the concentration of the conserved quantity, $\zeta \in L^{\infty}(\Omega)$ is the (possibly position-dependent) reaction coefficient and $q\in H^1(\Omega)'$ is the source. In thermal applications $\uP$ is the temperature, $\sP$ is the heat flux and $q-\zeta \uP$ is a heat source that depends linearly on the temperature. In weakly compressible flow in porous media $\uP$ is the pressure and $\sP$ is the fluid flux.

The boundary of $\Omega$, denoted by $\partial \Omega$, is divided into a Dirichlet part $\partial \Omega_D$, where $\uP=\uP_D$ is imposed, and a Neumann part $\partial\Omega_N$, in which $\sP \cdot \mathbf{\check{n}}=H$ is imposed. Hereafter we assume homogeneous boundary conditions throughout, i.e., $\uP_D = 0$ and $H = 0$.

Let us also define the field $\eP$ as the gradient of $\uP$,

\begin{equation}
	\eP = \nabla \uP
\label{eq:compatibility}
\end{equation}

This equation can be viewed as a compatibility condition. We call a vector field $\eP$ compatible if there exists a scalar field $\uP$ such that $\eP=\nabla \uP$.

Now, suppose that two vector fields $\stP$ and $\etP$ are given in $\Omega$. These fields will both be assumed in $\left[L^2(\Omega)\right]^{d_{\Omega}}$ and thus possibly discontinuous. Since they originate from a measured dataset, it is generally the case that $\stP$ does not satisfy the conservation law, and $\etP$ does not satisfy the compatibility condition. \textcolor{blue}{Moreover, they are concomitant, in the sense that each sample $\etP$ corresponds to the same spatial location as $\stP$ -- that is, they form coupled observations of gradient and flux at the same point in the domain --.}

\textcolor{red}{CGG: we have to make more clear that the fields $\etP$ and $\stP$ are concomitant.}

The goal, therefore, is to determine the best approximations to $\stP$ and $\etP$ within the space of feasible fields, that is, those that are both conservative and compatible. The scalar field $\uP$ emerges naturally as a by-product of this reconstruction.

\subsection{Well-posedness of the Exact Problem}

We begin by introducing the function spaces involved in the formulation:

\begin{equation}
	\left.
	\begin{array}{rcl}
		U & = & \left\{ \bar{u} \in H^1(\Omega),\ \bar{u} = 0 \text{ on } \partial\Omega_D \right\} \\
		E & = & \left[ L^2(\Omega) \right]^{d_\Omega} \\
		S & = & \left[ L^2(\Omega) \right]^{d_\Omega} \\
	\end{array}
	\right \rbrace  \quad X := U \times E \times S
\end{equation}

%\begin{eqnarray*} 
%	U&=&\{\bar{u}\in H^1(\Omega), \quad \bar{u}=0~\mbox{on}~\partial \Omega_D \}\\
%	E&=&\left[L^2(\Omega)\right]^{d_{\Omega}} \\
%	S&=&\left[L^2(\Omega)\right]^{d_{\Omega}} \\
%	X&=&U\times E \times S 
%\end{eqnarray*}


We now pose the problem as a minimization over the feasible set $ Z \subset X $, consisting of all tuples $ (\uP, \eP, \sP) \in X $ that satisfy the governing equations and boundary conditions.


\begin{tcolorbox}[colback=black!10!white, halign = justify, boxrule=0pt]
\begin{problem*}[P1]
	Find $x= \left (\uP,\eP,\sP \right) \in Z$ -- the feasible space -- that minimizes the distance to the data in the feasible space, i.e.,
	
	\begin{equation}
		J(x) \leq J(x+\delta x)~,\qquad \forall \,\delta x\,\in\,\delta Z~
	\end{equation}
	
	where 
	
	\begin{equation}
		J(\bar{\uP},\bar{\eP},\bar{s})=\frac12 \|\bar{\eP}-\etP\|^2 + \frac{\kappa}{2} \|\bar{\sP}-\stP\|^2~,
	\end{equation}
	
	and the feasible space is
	
	\begin{equation}
		Z=\left \{ (\bar{\uP},\bar{\eP},\bar{\sP})\in X~|~\nabla \cdot \bar{\sP} + \zeta \,\bar{\uP} = q~,~\bar{\eP}=\nabla \bar{\uP},\,\left.~\bar{\sP}\cdot\mathbf{\check{n}}\right |_{\partial \Omega_N}=0 \right \}
	\end{equation}
	
	with the variation given by
	
	\begin{equation}
		\delta Z=\left \{ (\delta{\uP},\delta{\eP},\delta{\sP})\in X~|~\nabla \cdot \delta{\sP} + \zeta \delta{\uP} = 0~,~ \delta{\eP}=\nabla \delta{\uP} ,\,\left. ~\delta{\sP}\cdot\mathbf{\check{n}}\right |_{\partial \Omega_N}=0\right \}~.
	\end{equation}
\end{problem*}
\end{tcolorbox}
\bigskip

\begin{remark}
	Above, $\|\cdot\|$ is the norm of $\left[L^2(\Omega)\right]^{d_{\Omega}}$. The corresponding scalar product is denoted by $\left ( \cdot , \cdot \right )$. The same notation will be used for $L^2(\Omega)$.
\end{remark}


Let us now relax the constraints through Lagrange multipliers $\lP$ and $\mP$, which belong to the following function spaces:

\begin{equation*}
	L=\{g\in H^1(\Omega)~|~g=0 \mbox{ on } \partial \Omega_D\},\qquad
	\mbox{ and }\qquad M=\left[L^2(\Omega)\right]^{d_{\Omega}},  
\end{equation*}

respectively. By the sake of simplicity, we will also use the following notation:

\begin{equation}
	y=\left (\lP,\mP\right )\,\in \, Y:=\,L\times M~.
\end{equation}

Then, the resulting equivalent problem can be written as: 

\begin{tcolorbox}[colback=black!10!white, halign = justify, boxrule=0pt]
	\begin{problem*}[P2]
		Find $(x,y)=\left( \uP, \eP, \sP, \lP, \mP \right) \subset X\times Y$ such that 
		\begin{equation}
			(x,y)= \mbox{arg} \left \{ \inf_{\bar{x}\in X} \sup_{\bar{y}\in Y} 
			\mathcal{L}(\bar{x},\bar{y})
			\right \}
		\end{equation}
		where
		\begin{equation}
			\begin{aligned}
				\mathcal{L}(\bar{x},\bar{y}) &~=~J(\bar{u},\bar{e},\bar{s})+ \nonumber \\
				&~-~\left ( \bar{s},\nabla \bar{\lambda} \right ) + \left (\zeta \bar{u} - q , \bar{\lambda} \right )  
				+ \left ( \bar{e}-\nabla \bar{u}, \bar{\mu} \right ) ~.
			\end{aligned} \label{eq:lagrangian}
		\end{equation}
	\end{problem*}
\end{tcolorbox}

The weak form of the optimality conditions reads:


\begin{equation}\label{eq:variational_form}
	\begin{array}{rcll}
		\left ( \lP, \zeta \delta \uP \right ) - \left ( \mP, \nabla \delta \uP \right ) & = & 0 &\qquad \forall \,\delta \uP \in U \\
		\left ( \eP, \delta \eP \right ) + \left ( \mP, \delta \eP \right ) &=& \left ( \etP, \delta \eP \right )  &\qquad \forall \,\delta \eP \in E\\
		\kappa \left ( \sP, \delta \sP \right ) - \left ( \nabla \lP, \delta \sP \right )  & = & \kappa \left ( \stP,\delta \sP \right ) &\qquad \forall\, \delta \sP \in S \\
		\left ( \zeta \,\uP, \delta \lP \right )-\left ( \sP, \nabla \delta \lP \right )  &=& \left ( q, \delta \lP \right ) 	&\qquad \forall\,\delta \lP \in L\\
		-\left (\nabla \uP,\delta \mP \right ) + \left ( \eP, \delta \mP \right ) & = & 0 &\qquad \forall \, \delta \mP \in M 
	\end{array}
\end{equation}

%\begin{eqnarray}
% \left ( \lambda, \zeta\delta u \right ) - \left ( \mu, \nabla \delta u \right ) & = & 0 \qquad\qquad\qquad\qquad \forall \,\delta u\in U  \label{eq:var1}\\
%\left ( e,\delta e \right )+ \left ( \mu, \delta e \right ) &=& \left ( \etP,\delta e \right )  \qquad\qquad\qquad \forall \,\delta e \in E\label{eq:var2}\\
%\kappa \left ( s, \delta s \right ) - \left ( \nabla \lambda, \delta s \right )  & = & \kappa \left ( \stP,\delta s \right ) \qquad\qquad\qquad \forall\, \delta s \in S \label{eq:var3}\\
%\left ( \zeta\,u,\delta \lambda \right )-\left ( s,\nabla \delta \lambda\right )  &=&\left ( q,\delta \lambda \right ) \qquad\qquad\qquad \forall\,\delta \lambda \in L\label{eq:var4}\\
%-\left (\nabla u,\delta \mu \right )+\left ( e,\delta \mu \right ) & = & 0 \qquad\qquad\qquad\qquad \forall\,\delta \mu \in M \label{eq:var5}
%\end{eqnarray}

The sum of the left-hand sides (with the last two equations multiplied by $-1$) defines the bilinear form of the problem in the product space $X \times Y$, given by:

\begin{equation}
	B((x,y),(\delta x,\delta y))=a(x,\delta x)+b(\delta x,y)-b(x,\delta y)
\end{equation}

where
\begin{equation}
	\begin{array}{rcl}
		a(x,\delta x)&=& \left (\eP,\delta \eP \right )+\kappa \left ( \sP,\delta \sP \right ) \\
		b(x,y)& = & \left ( \mP,\eP-\nabla \uP \right )-\left (\nabla \lP,\sP\right ) + \left ( \lP, \zeta \uP \right )
	\end{array}
\end{equation}

The optimality conditions in strong form, therefore, are given by:

\begin{eqnarray}
	\zeta \lP  + \nabla \cdot \mP & = & 0 			\label{eq:dif1}\\
	 \eP + \mP & = & \etP 							\label{eq:dif2}\\
	\kappa\, \sP - \nabla \lP & = & \kappa\,\stP 	\label{eq:dif3}\\
	 \zeta\,\uP+\nabla \cdot \sP  & = & q 			\label{eq:dif4}\\
	-\nabla \uP + \eP &=& 0 						\label{eq:dif5}
\end{eqnarray}

together with

\begin{eqnarray}
 \sP \cdot \mathbf{\check{n}} & = & 0~,\qquad \mbox{ on } \partial \Omega_N \\
 \mP \cdot \mathbf{\check{n}} & = & 0~,\qquad \mbox{ on } \partial \Omega_N ~.
\end{eqnarray}

\begin{remark}
	Notice that when the reaction coefficient $\zeta$ is identically zero the five variational equations given in (\ref{eq:variational_form}) -- or, equivalently, their differential versions, i.e., equations (\ref{eq:dif1}) to (\ref{eq:dif5}) -- are in fact two independent sets of equations: (a) The first, second and fifth equations, with unknowns $\uP$, $\eP$ and $\mP$, and (b) the third and fourth equations, with unknowns $\sP$ and $\lP$. In such a case, $\sP$ is the best approximation to $\stP$ among conservative flux fields, and $\eP=\nabla \uP$ is the best approximation to $\etP$ among compatible (gradient) fields, and both problems can be analyzed separately. The reaction term adds an interesting twist to the problem, coupling all five variational equations.	
\end{remark}

\begin{theorem}
	Problem P2 has a unique solution.
	
	\begin{proof}
		Following the theory of mixed problems \cite{brezzifortin,ernguermond} one needs to show that $a(\cdot,\cdot)$ and $b(\cdot,\cdot)$ are continuous, that $a(\cdot,\cdot)$ is coercive in $\delta Z$, with
		
		\begin{eqnarray}
			\delta Z&=&\{x\in X~|~b(x,y)=0\, \quad \forall y\in Y\}\nonumber \\
			&=&\{(\uP,\eP,\sP)\in X \mid \eP-\nabla \uP=0, ~\nabla \cdot \sP + \zeta \uP =0,~\left. ~\sP \cdot \mathbf{\check{n}}\right |_{\partial\Omega_N} = 0  \}~,
		\end{eqnarray}
	
		and that the inf-sup condition on $b(\cdot,\cdot)$,
		
		\begin{equation}
			\inf_{y\in Y}\sup_{x\in X} \frac{b(x,y)}{\|x\|_X\|y\|_Y} \geq K_b > 0~.
		\end{equation}
	
		holds true.
		
		The norms adopted in $X$ and $Y$ are 
		
		\begin{eqnarray}
			\|x\|_X &=& \|(\uP,\eP,\sP)\|_X = \sqrt{\|\nabla \uP\|^2+\|\eP\|^2+\kappa \|\sP\|^2}~,\\
			\|y\|_Y &=& \|(\lP,\mP)\|_Y=\sqrt{\frac{1}{\kappa}\left \| \nabla \lP \right \|^2+\left\|\mP \right \|^2}~.
		\end{eqnarray}
		
		\noindent \textcolor{red}{CGG: why not to consider $\frac{1}{\ell_\Omega}\|u\|^2$ and $\frac{1}{\ell_\Omega k }\|\lambda\|^2$ in the norm definition?}
		
		They were chosen so that they have the same units. It is easy to see that both $X$ and $Y$ are Hilbert spaces, and that $a(\cdot,\cdot)$ and $b(\cdot,\cdot)$ are continuous. Notice that
		
		\begin{equation*}
				a(x,x)=\|\eP\|^2 + \kappa \|\sP\|^2 ~.
		\end{equation*}

		So, for $x \in \delta Z$ it holds that:
		
		\begin{equation*}
			a(x,x)= \frac12 \|\eP\|^2 +\frac12 \|\nabla \uP\|^2+ \kappa \|\sP\|^2 \geq \frac12 \|x\|_X^2~.
		\end{equation*}

		Coercivity thus holds with constant equal to $\frac12$.
		
		\noindent \textcolor{red}{CGG: we can distribute with $c_1\|\nabla u \|^2 + (1-c_1) \|e\|^2$ and $c_1 \in (0,1)$.}
		
		It remains to show that there exists $K_b>0$ such that
		
		\begin{equation*}
			\sup_{(u,e,s)\in X} \frac{\left (\mP,\eP-\nabla \uP \right ) - \left (\nabla \lP, \sP \right ) + \left ( \lP, \zeta \uP \right )}{\|(\uP,\eP,\sP)\|_X} \geq K_b \, \|(\lP,\mP)\|_{Y}~, \qquad \forall (\lambda,\mu)\in Y~.
		\end{equation*}
	
		We can take $\uP=0$, $\eP=\mP$ and $\sP=-\frac{1}{\kappa}\nabla \lP$ to get
		
		\begin{equation*}
			\sup_{(\uP,\eP,\sP)\in X} \frac{\left (\mP,\eP-\nabla \uP \right ) - \left (\nabla \lP, \sP \right ) + \left ( \lP, \zeta\,\uP \right )}{\|(\uP,\eP,\sP)\|_X} 
			\geq \frac{\|\mP\|^2 + \frac{1}{\kappa}\|\nabla \lP\|^2}{\sqrt{\|\mP\|^2 + \frac{1}{\kappa}\|\nabla \lP \|^2}} = \|(\lP,\mP)\|_Y~,
		\end{equation*}

		giving an inf-sup constant $K_b$ equal to one. Thus, the coercivity of $a(\cdot,\cdot)$ on $\delta Z$ and the inf-sup stability of $b(\cdot,\cdot)$ have been proven. Then we have the well-posedness of the exact problem, with a unique solution $(\uP,\eP,\sP)\in H^1(\Omega) \times \left[L^2(\Omega)\right]^{d_{\Omega}} \times L^2(\Omega)$ that depends continuously on $q\in L^2(\Omega)$, $\etP \in \left[L^2(\Omega)\right]^{d_{\Omega}}$ and $\stP \in \left[L^2(\Omega)\right]^{d_{\Omega}}$.
	\end{proof}
\end{theorem}


\subsection{Approximation}

Let $U_h$, $E_h$, $S_h$, $L_h$ and $M_h$ be the finite element spaces for the five unknowns. We assume that $U_h\subset U$ and $L_h\subset L$ (in particular, the Dirichlet boundary conditions are strongly enforced). 
Also, we assume that $E_h$, $S_h$ and $M_h$ are subspaces of $L^2(\Omega)^d$ and define $X_h=U_h\times E_h \times S_h$ and $Y_h=L_h\times M_h$. The numerical approximation we propose corresponds to solving the following discrete problem:

\bigskip

\noindent{\bf Discrete Problem P2$_h$:} Find $(x_h,y_h)=\left ( u_h, e_h, s_h, \lambda_h, \mu_h\right ) \subset X_h\times Y_h$ such that 
\begin{equation}
(x_h,y_h)= \mbox{arg} \left \{ \inf_{\bar{x}\in X_h} \sup_{\bar{y}\in Y_h} 
\tilde{\mathcal{L}}(\bar{x},\bar{y})
\right \}
\end{equation}
where
\begin{eqnarray}
\tilde{\mathcal{L}}(\bar{x},\bar{y})&=&J(\bar{u},\bar{e},\bar{s})+ \nonumber \\
& &-\left ( \bar{s},\nabla \bar{\lambda} \right ) + \left (\zeta \bar{u} - q , \bar{\lambda} \right )  
+ \left ( \bar{e}-\nabla \bar{u}, \bar{\mu} \right ) +\nonumber \\
& &+\frac{\alpha}{2} \left \| \bar{e} - \nabla \bar{u} \right \|^2 + \\
& &-\frac{\gamma}{2} \|\bar{e}+\bar{\mu}-\etP\|^2 + \nonumber \\
& &-\frac{\eta}{2\kappa} \|\kappa \bar{s} - \nabla \bar{\lambda} - \kappa \stP\|^2 \nonumber \\
& &+ \frac{\theta \kappa \ell_s^2 }{2}\, \|\nabla \cdot \bar{s} + \zeta\, \bar{u} - q \|_h^2 + \nonumber \\
& & - \frac{\beta \ell_{\mu}^2 }{2}\, \|\nabla \cdot \bar{\mu} + \zeta\, \bar{\lambda} \|_h^2~.
\end{eqnarray}

\noindent \textcolor{red}{CGG: To keep the saddle-point structure, the squared norms of the optimality conditions for the dual variables are added and the squared norms of the optimality conditions for the primal variables are subtracted.}

\bigskip

The mesh-dependent norm $\|\cdot\|_h$ is defined as satisfying
\begin{equation}
\|f\|_h^2 = \sum_K \|f\|^2_{L^2(K)} 
\end{equation}
with the sum performed over all elements $K$ of the mesh. The associated scalar product is denoted as $\left ( \cdot, \cdot \right )_h$.

Above, $\alpha$, $\theta$ and $\beta$ are positive non-dimensional coefficients, and $\ell_s$ and $\ell_{\mu}$ are length scales to be determined later on. The function $\tilde{\mathcal{L}}$ is an augmented Lagrangian, of which the saddle point is the solution of P2$_h$. %%The solution is independent of $\alpha$ and $\theta$. 

\noindent \textcolor{red}{CGG: Up to here February 4, 2025.}


The optimality conditions of P2$_h$ are

\noindent \textcolor{red}{CGG: (37) fourth term - $\ell_s$ instead of $\ell_f$}

%\noindent \textcolor{red}{CGG: (38) third term - $(1+\gamma)$ instead of $(1-\gamma)$}


\begin{eqnarray}
\alpha \left ( \nabla u_h, \nabla \delta u \right ) + \theta \kappa \ell_s^2  \left ( \zeta^2 u_h, \delta u \right ) - \alpha \left ( e_h,\nabla \delta u \right ) +& & \nonumber \\
 + \theta \kappa \ell_f^2  \left ( \nabla \cdot s_h, \zeta \delta u \right )_h + \left ( \lambda_h, \zeta\delta u \right ) - \left ( \mu_h, \nabla \delta u \right ) & = & \theta \kappa \ell_s^2  \left ( q, \zeta\delta u \right )\nonumber \\ & & \quad \forall \,\delta u\in U_h  \label{eq:var1h}\\
-\alpha \left ( \nabla u_h, \delta e \right ) +(1+\alpha-\gamma)\,\left ( e_h,\delta e \right )+ (1-\gamma)\, \left ( \mu_h, \delta e \right ) &=& \left (1-\gamma)\,( \etP,\delta e \right )  \nonumber \\ & &\quad \forall \,\delta e \in E_h\label{eq:var2h}\\
(1-\eta)\kappa \left ( s_h, \delta s \right ) - (1-\eta)\left ( \nabla \lambda_h, \delta s \right ) + & & \nonumber \\
+\theta \kappa \ell_s^2  \left ( \nabla \cdot s_h, \nabla \cdot \delta s \right )_h+ \theta \kappa \ell_s^2 \,\left (\zeta\,u_h, \nabla \cdot \delta s \right )_h & = & (1-\eta)\kappa \left ( \stP,\delta s \right ) + \nonumber \\ & &+\theta \kappa \ell_s^2  \left ( q , \nabla \cdot \delta s \right )_h \nonumber \\ & & \quad \forall\, \delta s \in S_h \label{eq:var3h}\\
-(1-\eta)\left ( s_h,\nabla \delta \lambda\right ) + \left ( \zeta\,u_h,\delta \lambda \right ) - \beta \ell_{\mu}^2 \left ( \zeta^2 \lambda_h, \delta \lambda\right )+ & & \nonumber \\
- \beta \ell_{\mu}^2 \left ( \nabla \cdot \mu_h, \zeta\,\delta \lambda\right )_h -\frac{\eta}{\kappa} \left ( \nabla \lambda_h , \nabla \delta \lambda \right )&=&\left ( q,\delta \lambda \right ) + \eta \left ( \stP,\nabla \delta \lambda \right ) \nonumber \\ & & \quad \forall\,\delta \lambda \in L_h\label{eq:var4h}\\
-\left (\nabla u_h,\delta \mu \right )+(1-\gamma)\left ( e_h,\delta \mu \right )
- \beta \ell_{\mu}^2 \left ( \zeta\,\lambda_h, \nabla \cdot \delta \mu \right )_h + & & \nonumber \\- \beta \ell_{\mu}^2 \left ( \nabla \cdot \mu_h, \nabla \cdot \delta \mu \right )_h - \gamma \left ( \mu_h,\delta \mu \right ) & = & -\gamma \left (\etP,\delta \mu \right ) \nonumber \\ & & \quad \forall\,\delta \mu \in M_h \label{eq:var5h}
\end{eqnarray}

\bigskip

The sum of the left-hand sides (the last two multiplied by $-1$) is nothing but the discrete bilinear form that corresponds to P2$_h$, i.e.,
\begin{equation}
B_h((x_h,y_h),(\delta x,\delta y))=a_h(x_h,\delta x)+b(\delta x,y_h)-b(x_h,\delta y)+c_h(y_h,\delta y)~,
\end{equation}
where
\begin{eqnarray}
a_h(x_h,\delta x)&=&(1-\gamma)\left (e_h,\delta e\right )+(1-\eta)\kappa \left ( s_h,\delta s \right ) + \nonumber \\
& & +
\alpha \left ( e_h-\nabla u_h, \delta e - \nabla \delta u \right )+ \nonumber \\
& & + \theta \kappa \ell_s^2  \left ( \nabla \cdot s_h + \zeta u_h , \nabla \cdot \delta s + \zeta\,\delta u \right )_h~,
\end{eqnarray}
\begin{eqnarray}
b_h(\delta x,y_h)=b(\delta x,y_h)-\gamma \left (\mu_h,\delta e\right )+\eta \left (\nabla \lambda_h,\delta s \right ) ~,
\end{eqnarray}
\begin{eqnarray}
 c_h(y_h,\delta y)&= &  \beta \ell_{\mu}^2 \left ( \nabla \cdot \mu_h + \zeta \lambda_h , \nabla \cdot \delta \mu + \zeta\,\delta \lambda \right )_h+ \nonumber \\
 & & + \frac{\eta}{\kappa}\,\left ( \nabla \lambda_h,\nabla \delta \lambda \right ) + \gamma \left ( \mu_h,\delta \mu \right )~,
\end{eqnarray}
and $b(\cdot,\cdot)$ as defined in (\ref{eq:bform}).

Problem P2$_h$ can thus be written as "Find $(x_h,y_h)\in X_h \times Y_h$ such that
\begin{equation}
B_h\left ( \left ( x_h,y_h \right ),\left ( \delta x,\delta y \right ) \right )= F_h(\delta x,\delta y)
\label{eq:BzF}
\end{equation}
for all $(\delta x, \delta y)\in X_h\times Y_h$". The linear form $F_h:X_h\times Y_h \to \mathbb{R}$ stands for the sum of the right-hand sides of (\ref{eq:var1h})-(\ref{eq:var5h}), the last two multiplied by $-1$.

\bigskip

\noindent \textcolor{red}{CGG: Up to here February 5, 2025.}

\noindent{\bf Theorem 2:} Problem P2$_h$ is well-posed, and $(x_h,y_h)$ satisfies that $\exists c>0$ independent of $h$ such that
\begin{equation}
\||(x,y)-(x_h,y_h)\||_h \leq c\,\min_{(\bar{x},\bar{y})\in X_h\times Y_h}
\||(x,y)-(\bar{x},\bar{y})\||_h
\label{eq:bestapprox}
\end{equation}
where
\begin{eqnarray}
\||(x_h,y_h)\||_h^2 &=& \|\nabla u_h \|^2 +\|e_h\|^2 + \kappa \|s_h\|^2 + \kappa h^2 \|\nabla \cdot s_h\|_h^2 + \nonumber \\
& &+ \frac{1}{\kappa} \|\nabla \lambda_h\|^2 + \|\mu_h\|^2 + h^2\,\|\nabla \cdot \mu_h\|^2_h~.
\end{eqnarray}

\bigskip

The terms added to the Lagrangian allow us to work in a norm with improved control over $\nabla u_h$, $\nabla \cdot s_h$ and $\nabla \cdot \mu_h$.

\bigskip

\noindent{\em Proof:} It is easy to verify that the exact solution $(x,y)$ satisfies the discrete problem (\ref{eq:BzF}) when substituted in the place of $(x_h,y_h)$. The formulation is thus consistent.

 Let us establish the uniform strong coercivity of $B_h$, i.e., the existence of $k_{B}>0$ (independent of $h$) such that
\begin{equation}
B_h((x_h,y_h),(x_h,y_h))\geq k_B \left ( \|x_h\|_X^2 + \|y_h\|_Y^2 \right ),\qquad \forall (x_h,y_h)\in X_h\times Y_h~.
\end{equation}
By direct computation, 
\begin{eqnarray}
B_h((x_h,y_h),(x_h,y_h))&=&(1-\gamma) \|e_h\|^2 + (1-\eta)\kappa \|s_h\|^2 + \nonumber \\
& & + \alpha \|e_h-\nabla u_h\|^2 + \theta \kappa \ell_s^2 \|\nabla \cdot s_h + \zeta u_h\|^2 + \nonumber \\
& & + \beta \ell_{\mu}^2 \| \nabla \cdot \mu_h + \zeta \lambda_h \|^2
+ \frac{\eta}{\kappa} \|\nabla \lambda_h\|^2 + \gamma \|\mu_h\|^2~.
\end{eqnarray}


Now we use Young's inequality ($2ab\geq -\epsilon^2 a^2-(1/\epsilon^2)b^2$, $\forall a,b, \epsilon$) to each of the cross-product terms in the previous expression:

\begin{eqnarray}
-2\alpha \left (e_h,\nabla u_h \right ) & \geq & -\frac14 \|e_h\|^2 - 4\alpha^2 \|\nabla u_h \|^2 \\
%%%%%
2\theta\kappa \ell_s^2\left ( \nabla \cdot s_h, \zeta u_h \right )_h& \geq & - \frac12 \kappa \theta \ell_s^2 \|\nabla\cdot s_h\|^2 - {2\theta\kappa\ell_s^2}\|\zeta u_h\|^2\\
%%%%%%
2 \beta \ell_{\mu}^2\left ( \nabla \cdot \mu_h, \zeta \lambda_h \right )_h& \geq & - \frac12 {\beta} \ell_{\mu}^2 \|\nabla\cdot \mu_h\|^2 - {2\beta\ell_{\mu}^2}\|\zeta \lambda_h\|^2\\
%%%%%%%
\end{eqnarray}

Adding up all terms that contain $s_h$ we arrive at
\begin{equation}
(1-\eta) \kappa \|s_h\|^2 + \frac12 \kappa \ell_s^2 
 \theta  \|\nabla\cdot s_h\|_h^2~.
\label{eq:term1}
\end{equation}
Adding up all terms that contain $\zeta u_h$ gives
$$
-\theta \kappa \ell_s^2 \|\zeta u_h\|^2
$$
which using Poincar\'e-Friedrichs inequality
\begin{equation}
\|u_h\| \leq c_{\mbox{\scriptsize{PF}}}\,\ell_{\Omega}\,\|\nabla u_h\|, \qquad \forall u_h\in U_h
\label{eq:poincare}
\end{equation}
with $\ell_{\Omega}$ the diameter of $\Omega$ is shown to be greater or equal than 
\begin{equation}
- {\theta \kappa \ell_s^2 \|\zeta\|_{\infty}^2 c_{\mbox{\scriptsize{PF}}}^2 \ell_{\Omega}^2}\,\|\nabla u_h \|^2~.
\label{eq:term2}
\end{equation}
By $\|\zeta\|_{\infty}$ we denote the $L^{\infty}(\Omega)$-norm of $\zeta$.

Adding up all terms that contain $\mu_h$ we arrive at
\begin{equation}
\gamma \|\mu_h\|^2 +\frac12 \beta\ell_{\mu}^2  \|\nabla \cdot \mu_h \|_h^2~.
\label{eq:term4}
\end{equation}
The sum of the terms that contain $\zeta \lambda_h$ is greater or equal, by Poincaré-Friedrichs inequality, than
\begin{equation}
- \beta \ell_{\mu}^2 \|\zeta\|_{\infty}^2 c_{\mbox{\scriptsize{PF}}}^2 \ell_{\Omega}^2\,\|\nabla \lambda_h \|^2~.
\label{eq:term5}
\end{equation}
Adding up all terms that contain $\nabla u_h$ gives
\begin{equation}
\left ( \alpha - 4 \alpha^2 \right )
\|\nabla u_h\|^2~.
\label{eq:term6}
\end{equation}
Adding up all terms that contain $e_h$ gives
\begin{equation}
\left ( \frac34-\gamma + \alpha \right ) \|e_h\|^2~.
\label{eq:term7}
\end{equation}

Adding everything up we get
\begin{eqnarray}
B_h((x_h,y_h),(x_h,y_h))&\geq& \left (\frac34 -\gamma+\alpha \right ) \|e_h\|^2 + (1-\eta)\kappa \|s_h\|^2 + \nonumber \\
& & + (\alpha - 4 \alpha^2 - \theta \kappa \ell_s^2 \|\zeta\|_{\infty}^2 c_{\mbox{\scriptsize{PF}}}^2 \ell_{\Omega}^2)\, \|\nabla u_h\|^2 + \nonumber \\
& & + \frac12 \theta \kappa \ell_s^2 \|\nabla \cdot s_h\|^2 + \nonumber \\
& & + \frac12 \beta \ell_{\mu}^2 \| \nabla \cdot \mu_h\|^2 + \nonumber \\
& & 
+ \left ( \frac{\eta}{\kappa}-\beta \ell_{\mu}^2 \|\zeta\|_{\infty}^2 c_{\mbox{\scriptsize{PF}}}^2 \ell_{\Omega}^2 \right ) \|\nabla \lambda_h\|^2 + \gamma \|\mu_h\|^2~.
\end{eqnarray}

Taking
$$
\alpha = \frac{1}{8}, \quad \theta=\frac12,\quad \beta=\frac12,
\quad \gamma=\alpha, \quad \eta = \frac12,
\quad \ell_s = \ell_\mu = h, 
$$
we get
\begin{eqnarray*}
B_h\left ( (x_h,y_h),(x_h,y_h) \right ) &\geq  &
\left ( \frac{1}{16} - O(h^2) \right )\, \|\nabla u_h\|^2+\\
& & +\frac{3}{4}\,\|e_h\|^2+\frac12 \kappa \|s_h\|^2 + \frac14 \kappa h^2 \|\nabla \cdot s_h\|_h^2 + \\
& &
+\left (\frac{1}{2\kappa}  -O(h^2) \right ) \|\nabla \lambda_h\|^2 + \\
& & + \frac18 \|\mu_h\|^2  +\frac14 h^2 \|\nabla \cdot \mu_h\|_h^2 
\end{eqnarray*}
and finally
\begin{equation}
B_h\left ( (x_h,y_h),({x}_h,{y}_h) \right ) \geq 
k_B \, \| |(x_h,y_h)  \| |_h^2 
\end{equation}
where
\begin{equation}
k_B \geq \frac{1}{16}~.
\end{equation}
We have assumed that $h$ is sufficiently small so that the terms $O(h^2)$ can be neglected and thus $k_B$ becomes mesh independent. 

It is straightforward to see that $B_h(\cdot,\cdot)$ is continuous, i.e., $\exists N_B\geq 0$ such that
\begin{equation}
B_h\left ( (x_h,y_h),(\bar{x},\bar{y}) \right ) \leq N_B \||(x_h,y_h)\||_h\,\||(\bar{x},\bar{y})\||_h~,
\end{equation}
$ \forall (x_h,y_h) \in X_h\times Y_h$, $\forall (\bar{x},\bar{y})\in X_h\times Y_h$, with $N_B$ independent of $h$. Having proved that $B_h$ is consistent, continuous and coercive we obtain the best approximation property (\ref{eq:bestapprox}). $\Box$

\noindent \textcolor{red}{CGG: Up to here February 6, 2025.}



%=============================================
% --- Numerical experiments ---
\section{Numerical experiments}


In what follows we slightly change the notation and adopt bold symbols for the vector fields 
$\eP,~\sP$ and $\mP$.
In this section, we analyze the convergence properties of the discrete finite element approximation problem, considering the two formulations: (\textbf{I}) natural formulation, and (\textbf{II}) equal-order formulation. For the former, we consider 
\begin{equation}\label{eq:natural_func_spaces}
	\left(\uP_{h}, \eP_{h}, \sP_{h}, \lP_{h}, \mP_{h}\right) \in \CG^{D}_{k+1} \times \left[\DG_{k}\right]^{d_{\Omega}} \times \left[\DG_{k}\right]^{d_{\Omega}} \times \CG^{D}_{k+1} \times \left[\DG_{k}\right]^{d_{\Omega}} 
\end{equation} 
while, for the latter, we seek 
\begin{equation}\label{eq:equal_order_func_spaces}
\left(\uP_{h}, \eP_{h}, \sP_{h}, \lP_{h}, \mP_{h}\right)  \in \CG^{D}_{k+1} \times \left[\CG_{k+1}\right]^{d_{\Omega}} \times \left[\CG_{k+1}\right]^{d_{\Omega}} \times   \CG^{D}_{k+1} \times \left[\CG_{k+1}\right]^{d_{\Omega}}
\end{equation}

Here, the $k \in \mathbb{N}_{0}$ stands for the piecewise polynomial degree used to construct the finite element approximation space. The discrete solutions considered in this study were taken considering $k \in \{0, 1, 2\}$ for both formulations, which amounts to using linear, quadratic and cubic polynomial for the scalar fields $\uP$ and $\lP$. To assess the convergence properties of the method we consider both convex and non-convex geometries in two spatial dimensions. The convergence behavior in each case and for both formulations is evaluated through a manufactured solution chosen as a reference case. 

\subsection{Test case 1: a convex domain}

Let us first explain the procedure to construct a manufactured solution for 
this problem. First, consider the smooth field $\uP \in \mathcal{C}^{\infty}(\Omega)$ 
defined by:
\begin{equation}
	\uP:\Omega \to \mathbb{R},\quad \uP = \cos(\pi x)\cos(\pi y),
\end{equation}
whose gradient is given by:
\begin{equation*}
	\begin{aligned}
	\eP: \Omega \to \mathbb{R}^{2}, \quad \eP = [-\pi\sin(\pi x)\cos(\pi y), -\pi \cos(\pi x)\sin(\pi y)]^{\top}. \\
	\end{aligned}
\end{equation*}
Also define 
$$
\lP: \Omega \rightarrow \mathbb{R},\qquad \lP = -\frac{1}{40\,\pi}\,\left[\cos\left(2\,\pi\,x \right) + \cos\left(2\,\pi\,y \right)\right] \\ \\
$$
and
$$
\sP:\Omega \rightarrow \mathbb{R}^2, \qquad \sP = -\dfrac{\eP}{\|\eP\|}\, \left(\|\eP\| - \dfrac{1}{40}\,\|\eP\|^{3}\right)
$$
At this point, just for the purpose of assessment of the convergence properties of the discrete formulations 
under general conditions, it is convenient to 
introduce an additional source term to the right hand side of eq. (\ref{eq:dif1}), i.e.,
\begin{equation}
	\zeta \lP  + \nabla \cdot \mP  =  f
\end{equation}
By examining the optimality conditions in their strong form (\ref{eq:dif1})-(\ref{eq:dif5}) we
proceed as follows in order to construct the manufactured solution for the remaining fields.
First take
$$
\etP:\Omega \rightarrow \mathbb{R}^2, \qquad \etP = \eP + \frac{1}{20} \left[\sin\left(4\,\pi\,x \right), \sin\left(4\,\pi\,y\right) \right]^{\top}
$$
as being a perturbation of $\eP$, from which $\mP$ can be obtained using eq. (\ref{eq:dif2}).
Now, $\stP$ can be computed using eq. (\ref{eq:dif3}) and the source 
terms $f$ and $q$ are computed by substituting into eqs. (\ref{eq:dif1}) e (\ref{eq:dif4}). 
In the numerical experiments we take $\kappa = 1$ and $\zeta = 1$.

Notice that, this way, the fields $\eP$ and $\sP$ satisfy a sort of constitutive relation, whereas
the values realized by the fields $\etP$ and $\stP$ at the different points throughout the domain
does not necessarily do. This mimics practical scenarios in which $(\etP,\stP)$ are 
obtained from measured data.

Let $\Omega \subset \mathbb{R}^{2}$ such that $\Omega = [0,1]\times[0,1]$ whose boundary is defined as $\partial \Omega = \Gamma_{\text{left}} \cup \Gamma_{\text{right}} \cup \Gamma_{\text{bottom}} \cup \Gamma_{\text{top}}$.
The problem is solved imposing the following boundary conditions
\begin{equation}
	\left\lbrace
	\begin{array}{rclll}
		\sP \cdot \mathbf{\check{n}} 			&=& \mathsf{g}_{\sP} 			&\mbox{on}&~\Gamma_{\text{bottom}} \cup \Gamma_{\text{top}} \\
		\mP \cdot \mathbf{\check{n}} &=& \mathsf{g}_{\mP} &\mbox{on}&~\Gamma_{\text{bottom}} \cup \Gamma_{\text{top}} \\
		\uP &=& \mathsf{g}_{\uP} &\mbox{on}&~\Gamma_{\text{left}} \cup \Gamma_{\text{right}} \\
		\lP &=& \mathsf{g}_{\lP} &\mbox{on}&~\Gamma_{\text{left}} \cup \Gamma_{\text{right}} \\
	\end{array}
	\right.
\end{equation}
where the values in the right hand side are computed from the manufactured analytical fields.
In order to illustrate this manufactured solution, Figure~\ref{fig:all_fields_ref}(a) shows
a plot of the norm of $\sP$ as a function of the norm of $\eP$, along with their counteparts 
$\stP$ and $\etP$ evaluated over a computational mesh. 
Notice that the exact flux-gradient relationship is depicted as a strictly monotonic curve (thick black
line), whereas the fields $\stP$ and $\etP$ (blue dots) exhibit a certain distribution around the 
former.
\begin{figure}[h!]
	\centering
	\caption{Figure (a) presents the functional relationship between flux and gradient norms for both exact and perturbed cases: $\|s\|~\mbox{vs.}~\|e\|$, and $\|\stP\|~\mbox{vs.}~\|\etP\|$, respectively. Figure (b) to (f) presents the $L^{2}(\Omega)$-orthogonal projections of fields $u$, $\lP$, $e$, $s$, and $\mP$ onto the corresponding function spaces defined by the natural formulation, respectively.}
	%\fontsize{11pt}{10pt}\selectfont 
	\def\svgwidth{1.0\columnwidth}
	\input{Figures/PDFs/all_fields_ref.pdf_tex}
	\label{fig:all_fields_ref}
\end{figure}
We consider a sequence of finite element conformal and shape-regular partitions of the computational domain $\mathcal{T}_h$ with characteristic mesh size $h$. 

We first examine the natural formulation (\textbf{I}) with no stabilization. By setting $ k = 0 $ in eq. \ref{eq:natural_func_spaces}, the scalar fields satisfy $ \uP_{h}, \lP_{h} \in \CG_{1}^{D} \subset H^{1}(\Omega) $.  


Theorem~2 ensures that the convergence rate of the global error is governed by the dominant term in the mesh-dependent norm. In this case, the dominant contribution is of order $ \mathcal{O}(h) $, as the errors in $ \nabla \uP_h $ and $ \nabla \lP_h $ scale as $ \mathcal{O}(h) $, whereas the errors in $ \eP_h $, $ \sP_h $, and $\mP_h $ scale as $ \mathcal{O}(h^2) $. The numerical results shown in Figures~\ref{fig:convergence_scalar_linear}(a) and~\ref{fig:convergence_scalar_linear}(b) are in agreement with these theoretical error estimates.

Similarly, Figures ~\ref{fig:convergence_scalar_linear}(c)--\ref{fig:convergence_scalar_linear}(d) and Figures ~\ref{fig:convergence_scalar_linear}(e)--\ref{fig:convergence_scalar_linear}(f) present the convergence behavior of $ \uP_{h} $ and $ \lP_{h} $ in the equal-order formulation, respectively, and under the same configuration employed previously. For both fields, the $ L^{2} $- and $ H^{1} $- norms were analyzed. In this case, the convergence rates are identical to those observed in the natural formulation. This behavior is in agreement with theory since these fields are approximated by the same space of discrete functions in both formulations. Furthermore, no differences were observed between the convergence rates of the unstabilized, minimally stabilized, or fully stabilized formulations for $\uP_{h}$ and $\lP_{h}$ in both standards, although the fully stabilized formulation presented a better error behavior for the $H^{1}$-norm for $\lP_{h}$, as shown in Figure \ref{fig:convergence_scalar_linear}(f).


Analyzing the vector fields under identical choices of $k$ and $\zeta$, Figure~\ref{fig:convergence_vector_linear}(a) depicts the convergence of all vector quantities, i.e. $\sP_{h}$, $\eP_{h}$, and $\mP_{h}$, in $L^{2}$-norm for the natural formulation. In this case, linear convergence was observed for all the fields. Figures~\ref{fig:convergence_vector_linear}(b)-\ref{fig:convergence_vector_linear}(c), on the other hand, presents the convergence of the vector unknowns in $L^{2}$-norm for the equal-order formulation. In this case, the flux vector field, $\sP_{h}$, exhibited identical convergence behavior across unstabilized, minimally stabilized, and fully stabilized formulations, presenting overlapping convergence curves likewise for the scalar fields. 

For $\eP_{h}$ and $\mP_{h}$, however, the minimally and fully stabilized curves presented did not reach a second-order convergence rate, like the unstabilized one. Although both do not violate the expected theoretical error estimate, we should expect that the introduction of the stabilization parameter is linked to a gain in the convergence rate or a reduction in the error magnitude.

Consequently, analogous numerical experiments were performed employing $\text{P}2$-elements ($k=1$). For scalar fields  and  in both norms within the natural formulation, convergence improved substantially, achieving third-order convergence in the $L^{2}$-norm and second-order convergence in the $H^{1}$-norm, as illustrated in Figures~\ref{fig:convergence_scalar_quadratic}(a)--\ref{fig:convergence_scalar_quadratic}(b). Under the equal-order formulation, the introduction of stabilization parameters notably enhanced convergence, with minimally and fully stabilized formulations achieving third-order convergence rates in the $L^{2}$-norm and second-order rates in the $H^{1}$-norm, as depicted in Figures~\ref{fig:convergence_scalar_quadratic}(c)--\ref{fig:convergence_scalar_quadratic}(f). Conversely, the unstabilized formulation exhibited no gain from polynomial refinement, maintaining the convergence rates seen with $\text{P}1$-elements.


Examining the vector quantities with $\text{P}2$-elements in the natural formulation (Figure~\ref{fig:convergence_vector_quadratic}(a)), convergence in the $L^{2}$-norm followed the anticipated second-order rate. However, in the equal-order formulation, the effectiveness of stabilization parameters became evident. Specifically, a third-order convergence rate for the flux field $\sP_{h}$ was observed solely with the fully stabilized formulation (Figure~\ref{fig:convergence_vector_quadratic}(b)). This aligns with theoretical expectations since the parameter $\theta$, exclusive to the fully stabilized formulation, explicitly contributes to controlling the flux divergence.

For fields $\eP_{h}$ and $\mP_{h}$, however, polynomial refinement provided no further rate improvement, with second-order convergence consistently observed (Figures~\ref{fig:convergence_vector_quadratic}(c)--\ref{fig:convergence_vector_quadratic}(d)). Nonetheless, the error magnitudes were significantly reduced under the fully stabilized formulation compared to other cases.

In addition, numerical experiments employing $\text{P}3$-elements ($k=2$) demonstrated further improvements in convergence rates. Under the natural formulation, scalar and vector fields showed fourth-order and third-order convergence rates in the $L^{2}$- and $ H^{1} $-norms for $\uP_{h}$ and $ \lP_{h} $, respectively. Additionally, a third-order convergence rate was observed in the $L^{2}$-norm for $\sP_{h}$, $\eP_{h}$, and $\mP_{h}$.

Under the equal-order formulation, the minimally and fully stabilized cases achieved superior convergence rates for scalar fields, attaining fourth-order convergence in the $L^{2}$-norm and third-order convergence in the $H^{1}$-norm, in contrast to the unstabilized formulation, which reached at most a third-order rate. For vector fields, the fully stabilized formulation for $\sP_{h}$ achieved a fourth-order convergence rate, while the minimally stabilized formulation exhibited intermediate behavior, and the unstabilized case retained third-order convergence. The fields $\eP_{h}$ and $\mP_{h}$ displayed trends similar to those observed with $\text{P}2$-elements, with notably improved convergence under full stabilization.

Convergence rates in the $ \text{H(div)} $-norm were also investigated for the vector fields under the equal-order formulation across all three polynomial degrees. Although theoretical error estimates for this norm are not available, the fully stabilized formulation consistently produced the most favorable results in all scenarios considered in this study. For $ \text{P}1 $-elements, no significant differences were observed in the convergence behavior of the flux $ \sP_{h} $ across the unstabilized, minimally stabilized, and fully stabilized formulations, with the corresponding error curves largely overlapping. In contrast, for $ \eP_{h} $ and $ \mP_{h} $, the fully stabilized formulation exhibited sub-linear convergence, unlike the unstabilized case. However, it achieved lower error magnitudes overall.

For $ \text{P}2 $-elements, the fully stabilized formulation yielded the best performance across all vector variables. In particular, it was the only one to attain second-order convergence for the flux $ \sP_{h} $, and it also resulted in the smallest error magnitudes for $ \eP_{h} $ and $ \mP_{h} $. A similar trend was observed for $ \text{P}3 $-elements. The enhanced convergence of the flux field in the $ \text{H(div)} $-norm for higher-order elements under the fully stabilized formulation can be attributed to the presence of the parameter $ \theta $, which explicitly controls the divergence of the flux in the discrete norm.

To conclude this section, we emphasize that, for all fields, additional numerical experiments were carried out to evaluate the effect of each individual stabilization parameter in isolation, in comparison with the three primary cases: unstabilized, minimally stabilized, and fully stabilized. The error curves corresponding to the isolated parameters displayed only minor variations in magnitude and exhibited convergence rates consistent with those of the main reference cases. Furthermore, numerical experiments confirmed that the convergence behavior is largely independent of the reaction coefficient $ \zeta $, as similar results were obtained when varying $ \zeta $ over several orders of magnitude, from $ 10^{-3} $ to $ 10^{3} $.

%=============================================
% k = 0

\begin{figure}[h!]
	\centering
	\caption{Figure (a) and (b) show the convergence of the scalar quantities $\uP_{h}$ and $\lP_{h}$ in the natural formulation with $k=0$, measured in the $\|\cdot\|_{L^{2}}$ and $\|\cdot\|_{H^{1}}$ norms, respectively. Figures (c) and (d) illustrate the convergence of the field $u$ in the equal-order formulation with $k=0$, evaluated in the same norms. The results are reported for three stabilization configurations: unstabilized, minimally stabilized ($ \alpha, \eta $), and fully stabilized. Similarly, Figures~(e) and (f) present the convergence of $ \lP_{h} $ in the equal-order formulation under the same stabilization settings and norms.}
	%\fontsize{11pt}{10pt}\selectfont 
	\def\svgwidth{1.0\columnwidth}
	\input{Figures/PDFs/convergence_scalar_linear.pdf_tex}
	\label{fig:convergence_scalar_linear}
\end{figure}


\begin{figure}[h!]
	\centering
	\caption{Figure (a) show the convergence of the vector quantities $\sP_{h}$, $\eP_{h}$, and $\mP_{h}$ in the natural formulation with $k=0$, measured in the $\|\cdot\|_{L^{2}}$ norm, respectively. Figures (b), (c) and (d) illustrate the convergence of the fields $\sP_{h}$, $\eP_{h}$, and $\mP_{h}$ in the equal-order formulation with $k=0$, also in the $\|\cdot\|_{L^{2}}$. The results are presented for three stabilization configurations: the unstabilized, the fully stabilized, and the minimally stabilzed $(\alpha,\eta)$.}
	%\fontsize{11pt}{10pt}\selectfont 
	\def\svgwidth{1.0\columnwidth}
	\input{Figures/PDFs/convergence_vector_linear.pdf_tex}
	\label{fig:convergence_vector_linear}
\end{figure}


\begin{figure}[h!]
	\centering
	\caption{Figure (a) to (c) show the convergence of the vector quantities $\sP_{h}$, $\eP_{h}$, and $\mP_{h}$ in the equal-order formulation with $k=0$, measured in the $\|\cdot\|_{\text{Hdiv}}$ norm, respectively. For all of them, three stabilization configurations are considered: the unstabilized, the fully stabilized, and the and the minimally stabilized $(\alpha, \eta)$.}
	%\fontsize{11pt}{10pt}\selectfont 
	\def\svgwidth{1.0\columnwidth}
	\input{Figures/PDFs/convergence_vector_linear_hdiv.pdf_tex}
	\label{fig:convergence_vector_linear_hdiv}
\end{figure}



%=============================================
% k = 1

\begin{figure}[h!]
	\centering
	\caption{Figure (a) and (b) show the convergence of the scalar quantities $u$ and $\lP_{h}$ in the natural formulation with $k=1$, measured in the $\|\cdot\|_{L^{2}}$ and $\|\cdot\|_{H^{1}}$ norms, respectively. Figures (c) and (d) illustrate the convergence of the field $\uP_{h}$ in the equal-order formulation with $k=1$, also in the $\|\cdot\|_{L^{2}}$ and $\|\cdot\|_{H^{1}}$. The latter considers three cases: the unstabilized, the fully stabilized, and the minimally stabilized $(\alpha, \eta)$. Similarly, Figures (e) and (f) present the convergence of $\lP_{h}$ in the equal-order formulation under the same stabilization configurations and norms as used for $\uP_{h}$.}
	%\fontsize{11pt}{10pt}\selectfont 
	\def\svgwidth{1.0\columnwidth}
	\input{Figures/PDFs/convergence_scalar_quadratic.pdf_tex}
	\label{fig:convergence_scalar_quadratic}
\end{figure}



\begin{figure}[h!]
	\centering
	\caption{Figure (a) show the convergence of the vector quantities $\sP_{h}$, $\eP_{h}$, and $\mP_{h}$ in the natural formulation with $k=1$, measured in the $\|\cdot\|_{L^{2}}$ norm, respectively. Figures (b), (c) and (d) illustrate the convergence of the fields $\sP_{h}$, $\eP_{h}$, and $\mP_{h}$ in the equal-order formulation with $k=1$, also in the $\|\cdot\|_{L^{2}}$. The latter considers three cases: the unstabilized, the fully stabilized, and the and the minimally stabilized $(\alpha, \eta)$.}
	%\fontsize{11pt}{10pt}\selectfont 
	\def\svgwidth{1.0\columnwidth}
	\input{Figures/PDFs/convergence_vector_quadratic.pdf_tex}
	\label{fig:convergence_vector_quadratic}
\end{figure}

\begin{figure}[h!]
	\centering
	\caption{Figure (a) to (c) show the convergence of the vector quantities $\sP_{h}$, $\eP_{h}$, and $\mP_{h}$ in the equal-order formulation with $k=1$, measured in the $\|\cdot\|_{\text{Hdiv}}$ norm, respectively. For all of them, three cases are considered: the unstabilized, the fully stabilized, and the and the minimally stabilized $(\alpha, \eta)$.}
	%\fontsize{11pt}{10pt}\selectfont 
	\def\svgwidth{1.0\columnwidth}
	\input{Figures/PDFs/convergence_vector_quadratic_hdiv.pdf_tex}
	\label{fig:convergence_vector_quadratic_hdiv}
\end{figure}




%=============================================
% k = 2
\begin{figure}[h!]
	\centering
	\caption{Figure (a) and (b) show the convergence of the scalar quantities $\uP_{h}$ and $\lP_{h}$ in the natural formulation with $k=2$, measured in the $\|\cdot\|_{L^{2}}$ and $\|\cdot\|_{H^{1}}$ norms, respectively. Figures (c) and (d) illustrate the convergence of the field $\uP_{h}$ in the equal-order formulation with $k=2$, also in the $\|\cdot\|_{L^{2}}$ and $\|\cdot\|_{H^{1}}$. The latter considers three cases: the unstabilized, the fully stabilized, and the minimally stabilized $(\alpha, \eta)$. Similarly, Figures (e) and (f) present the convergence of $\lP_{h}$ in the equal-order formulation under the same stabilization configurations and norms as used for $\uP_{h}$.}
	%\fontsize{11pt}{10pt}\selectfont 
	\def\svgwidth{1.0\columnwidth}
	\input{Figures/PDFs/convergence_scalar_cubic.pdf_tex}
	\label{fig:convergence_scalar_cubic}
\end{figure}



\begin{figure}[h!]
	\centering
	\caption{Figure (a) show the convergence of the vector quantities $\sP_{h}$, $\eP_{h}$, and $\mP_{h}$ in the natural formulation with $k=2$, measured in the $\|\cdot\|_{L^{2}}$ norm, respectively. Figures (b), (c) and (d) illustrate the convergence of the fields $\sP_{h}$, $\eP_{h}$, and $\mP_{h}$ in the equal-order formulation with $k=2$, also in the $\|\cdot\|_{L^{2}}$. The latter considers three cases: the unstabilized, the fully stabilized, and the and the minimally stabilized $(\alpha, \eta)$.}
	%\fontsize{11pt}{10pt}\selectfont 
	\def\svgwidth{1.0\columnwidth}
	\input{Figures/PDFs/convergence_vector_cubic.pdf_tex}
	\label{fig:convergence_vector_cubic}
\end{figure}

\begin{figure}[h!]
	\centering
	\caption{Figure (a) to (c) show the convergence of the vector quantities $\sP_{h}$, $\eP_{h}$, and $\mP_{h}$ in the equal-order formulation with $k=2$, measured in the $\|\cdot\|_{\text{Hdiv}}$ norm, respectively. For all of them, three cases are considered: the unstabilized, the fully stabilized, and the and the minimally stabilized $(\alpha, \eta)$.}
	%\fontsize{11pt}{10pt}\selectfont 
	\def\svgwidth{1.0\columnwidth}
	\input{Figures/PDFs/convergence_vector_cubic_hdiv.pdf_tex}
	\label{fig:convergence_vector_cubic_hdiv}
\end{figure}


\cleardoublepage
\newpage
\subsection{Numerical convergence under non-convex domain}

In continuation of the previous approach, we now consider the case of a non-convex domain. So, let $\Omega \subset \mathbb{R}^2$ be a domain with a reentrant corner at the origin $(0, 0)$ and an internal angle of $\phi \in \left\{\frac{\pi}{4}, \frac{\pi}{2}, \frac{3\pi}{4}\right\}$. Formally, the computational domain can be defined as: $\Omega = [-1,1] \times [-1,1] \setminus C_{\phi}$, where $C_\phi = \left\{ (x, y) \in \mathbb{R}^2 \;\middle|\; 0 \le x \le 1,\; -\tan(\phi)x \le y \le 0 \right\}$. Figure \ref{fig:geometries_non_convex} illustrate the three different domains under consideration. Here, consider the scalar field $u$ as being:

\begin{equation*}
	\begin{aligned}
	&u: \Omega \rightarrow \mathbb{R}, \quad u = r^{\nu}\,\sin(\nu\,\theta) \\
	&\mbox{where}~r = \sqrt{\left(x^{2}+y^{2}\right)}, \quad \theta = \arctan{\left(\frac{y}{x}\right)}, \quad  \nu = \frac{\pi}{\psi},~\mbox{with}~\psi = 2\pi-\phi
	\end{aligned}
\end{equation*}

The corresponding gradient field is then given explicity by

\begin{equation*}
	\begin{aligned}
	e: \Omega \rightarrow \mathbb{R}^{2}, \quad e = \nu\,r^{\nu-2}\left[x\,\sin(\nu\,\theta)-y\,\cos(\nu\,\theta),y\,\sin(\nu\,\theta)+x\,\cos(\nu\,\theta)\right]^{\top}
	\end{aligned}
\end{equation*}

In this setting, the expressions for $\lP:\Omega \rightarrow \mathbb{R}$, $\etP:\Omega \rightarrow \mathbb{R}^{2}$, and $s:\Omega \rightarrow \mathbb{R}^{2}$ are taken to be identical to those employed in the preceding numerical experiment, with the exception that the coefficient $A_{s}$ is set to zero in the definition of the exact flux. inally, it is important to emphasize that the current problem assumes pure Dirichlet boundary conditions, in contrast to the mixed boundary conditions considered previously.

\begin{figure}[h!]
	\centering
	\caption{Illustration of the three non-convex geometries considered.}
	\fontsize{11pt}{10pt}\selectfont 
	\def\svgwidth{1.0\columnwidth}
	\input{Figures/PDFs/geometries_non_convex.pdf_tex}
	\label{fig:geometries_non_convex}
\end{figure}

To illustrate the manufactured field we are dealing with, Figure~\ref{fig:all_fields_ref_non_convex}(a) presents the functional relationship between the norms of the exact flux and the exact gradient ($ \sP, \eP $), along with the corresponding relationship for their perturbed counterparts ($ \stP, \etP $).  It is worth noting that, in this case, the cubic term has been suppressed from the constitutive law. Furthermore, Figures~\ref{fig:all_fields_ref_non_convex}(b)--\ref{fig:all_fields_ref_non_convex}(f) present the individual manufactured fields over the domain shown in Figure~\ref{fig:geometries_non_convex}(a), which corresponds to the acute angle configuration.

\begin{figure}[h!]
	\centering
	\caption{Figure (a) presents the functional relationship between flux and gradient norms for both exact and perturbed cases: $\|\sP\|~\mbox{vs.}~\|\eP\|$, and $\|\stP\|~\mbox{vs.}~\|\etP\|$, respectively. Figure (b) to (f) presents the $L^{2}(\Omega)$-orthogonal projections of fields $\uP$, $\lP$, $\eP$, $\sP$, and $\mP$ onto the corresponding function spaces defined by the natural formulation, respectively.}
	%\fontsize{11pt}{10pt}\selectfont 
	\def\svgwidth{1.0\columnwidth}
	\input{Figures/PDFs/all_fields_ref_non_convex.pdf_tex}
	\label{fig:all_fields_ref_non_convex}
\end{figure}


For all these geometries, it can be shown that $\uP$ possesses fewer than $\nu+1$ derivatives in $L^{2}$, implying that $\uP \in H^{1}(\Omega)$, but $\uP \notin H^{2}(\Omega)$. Consequently, no improvement in the convergence rate is expected from polynomial function space refinement. To overcome this limitation, a local mesh refinement near the singularity must be employed, as illustrated, for instance, in Figure~\ref{fig:geometries_non_convex}(a).

So, instead of varying the polynomial degree of the discrete function spaces, we opt to use $\text{P}1$-elements ($ k=0 $) and we assess the convergence behavior with respect to different corner angles. Specifically, for the values $ \phi \in \left\{ \frac{\pi}{4}, \frac{\pi}{2}, \frac{3\pi}{4} \right\} $, the corresponding values of $ \psi $ are $ \left\{ \frac{7\pi}{4}, \frac{3\pi}{2}, \frac{5\pi}{4} \right\} $. Consequently, the regularity of the scalar field $ \uP $ is limited to the spaces $ H^{11/7 - \varepsilon}(\Omega) $, $ H^{5/3 - \varepsilon}(\Omega) $, and $ H^{9/5 - \varepsilon}(\Omega) $, for all $ \varepsilon > 0 $. Error estimates for finite element schemes in domains with corners have been extensively studied in the literature, for instance by \cite{Babuska1972, Babuska1970}. As a result, the convergence rate of $ \uP $ in both $ L^{2} $-norm  and $ H^{1} $-norm are expected to be restricted by the singular behavior induced by the reentrant corner and one may anticipate convergence rates of  $\mathcal{O}(h^{1+4/7-\varepsilon})$, $ \mathcal{O}(h^{1+2/3-\varepsilon}) $ and $ \mathcal{O}(h^{1+4/5-\varepsilon}) $ for all $ \varepsilon > 0 $ in $L^{2}$-norm, and $\mathcal{O}(h^{4/7-\varepsilon})$, $ \mathcal{O}(h^{2/3-\varepsilon}) $ and $ \mathcal{O}(h^{4/5-\varepsilon}) $ for all $ \varepsilon > 0 $ in $H^{1}$-norm. 

Table~\ref{tab:scalar_errors} summarizes the convergence results for the scalar fields in both $ L^{2} $- and $ H^{1} $-norms, considering the natural formulation, the equal-order formulation, and its stabilized variants, across the three corner angles. Similarly, Table~\ref{tab:vector_errors} reports the convergence behavior of the vector fields in the $ L^{2} $-norm under the same formulations and geometric configurations.


Focusing on the scalar field $ \uP_h $ under the natural formulation, we observe that the convergence in the $ H^1 $-norm closely approaches the theoretical upper bounds dictated by the corner singularity. In the $ L^2 $-norm, the convergence rate approximately doubles that of the $ H^1 $-norm, as expected from standard interpolation theory. The field $ \lP_h $ exhibits slightly better convergence than $ \uP_h $, particularly in the $ H^1 $-norm.

\begin{table}[h!]
	\centering
	\caption{Estimated convergence rates for the scalar fields $ \uP_h $ and $ \lP_h $ in the $ L^2 $- and $ H^1 $-norms, under different stabilization schemes and reentrant corner angles $ \psi \in \left\{ \frac{7\pi}{4}, \frac{3\pi}{2}, \frac{5\pi}{4} \right\} $.}
	%\rowcolors{3}{black!10!white}{white}
	\begin{tabular}{l l l l l}
		\toprule
		%\rowcolor{black!30!white} 
		Formulation & Norms & $\psi = \frac{7\pi}{4}$ & $\psi = \frac{3\pi}{2}$ & $\psi = \frac{5\pi}{4}$ \\ [2pt]
		\hline
		\multirow{4}{*}{Natural}				 & $\|u-\uP_{h}\|_{L^{2}}$ & 1.14  & 1.44  & 1.65  \\
		                                 & $\|u-\uP_{h}\|_{H^{1}}$ & 0.57 & 0.72  & 0.81 \\
		                                 & $\|\lP-\lP_{h}\|_{L^{2}}$ & 1.22  & 1.68 & 1.97  \\
		                                 & $\|\lP-\lP_{h}\|_{H^{1}}$ & 0.99  & 0.99 & 0.99  \\
		\hline
		\multirow{4}{*}{Unstabilized}    & $\|u-\uP_{h}\|_{L^{2}}$ & 1.43 & 1.65 & 1.61 \\
		                                 & $\|u-\uP_{h}\|_{H^{1}}$ & 0.46 & 0.59 & 0.64 \\
		                                 & $\|\lP-\lP_{h}\|_{L^{2}}$ & 1.16 & 1.30 & 2.16 \\
		                                 & $\|\lP-\lP_{h}\|_{H^{1}}$ & 0.46 & 0.30 & 1.24 \\
		\hline
		\multirow{4}{*}{Fully stabilized} & $\|u-\uP_{h}\|_{L^{2}}$ & 1.20 & 1.64 & 1.84 \\
		                                 & $\|u-\uP_{h}\|_{H^{1}}$ & 0.57 & 0.72 & 0.81 \\
		                                 & $\|\lP-\lP_{h}\|_{L^{2}}$ & 1.10 & 1.39 & 1.82 \\
		                                 & $\|\lP-\lP_{h}\|_{H^{1}}$ & 0.60 & 0.82 & 0.99 \\
		\hline
		\multirow{4}{*}{Minimally stabilized} & $\|u-\uP_{h}\|_{L^{2}}$ & 1.16 & 1.51 & 1.82 \\
		                                     & $\|u-\uP_{h}\|_{H^{1}}$ & 0.57 & 0.73 & 0.81 \\
		                                     & $\|\lP-\lP_{h}\|_{L^{2}}$ & 1.11 & 1.48 & 1.93 \\
		                                     & $\|\lP-\lP_{h}\|_{H^{1}}$ & 0.60 & 0.78 & 0.98 \\
		\bottomrule
	\end{tabular}
	\label{tab:scalar_errors}
\end{table}

When comparing the same fields under the equal-order formulation, the convergence rates remain similar in magnitude to those obtained with the natural formulation, with no substantial improvement. Among the equal-order variants, we highlight the enhanced convergence of both $ \uP_h $ and $ \lP_h $ in the $ H^1 $-norm achieved by the minimally and fully stabilized formulations, compared to the unstabilized case. However, no further improvements are observed in the $ L^2 $-norm for either field when stabilization is applied.

For the vector fields $ \sP_h $, $ \eP_h $, and $ \mP_h $, we observe that the convergence rates in the $ L^2 $-norm remain nearly identical across all formulations for each fixed corner angle. Notable improvements in convergence are primarily associated with increasing the interior angle at the reentrant corner, with higher accuracy consistently achieved when transitioning from acute to obtuse geometries. 

These results confirm that the convergence behavior is primarily dictated by the corner-induced singularity. While stabilization enhances robustness -- especially in the $H^1$-norm for scalar fields -- it does not overcome the inherent regularity limitations imposed by the geometry. 


\begin{table}[h!]
	\centering
	\caption{Estimated convergence rates in the $ L^2 $-norm for the vector fields $ \sP_h $, $ \eP_h $, and $ \mP_h $, under various stabilization schemes and for different reentrant corner angles $ \psi \in \left\{ \frac{7\pi}{4}, \frac{3\pi}{2}, \frac{5\pi}{4} \right\} $.}
	%\rowcolors{3}{black!10!white}{white}
	\begin{tabular}{l l c c c}
		\toprule
		%\rowcolor{black!30!white} 
		Formulation & Norms & $\psi = \frac{7\pi}{4}$ & $\psi = \frac{3\pi}{2}$ & $\psi = \frac{5\pi}{4}$ \\ [2pt]
		\hline
		\multirow{3}{*}{Natural}    & $\|\mathbf{s}-\sP_{h}\|_{L^{2}}$ & 0.57 & 0.73 & 0.83  \\
		                                 & $\|\mathbf{e}-\eP_{h}\|_{L^{2}}$ & 0.57 & 0.72 & 0.81 \\
		                                 & $\|\mP-\mP_{h}\|_{L^{2}}$ & 0.58 & 0.73 & 0.87 \\
		\hline
		\multirow{3}{*}{Unstabilized}    & $\|\mathbf{s}-\sP_{h}\|_{L^{2}}$ & 0.55 & 0.67 & 0.86 \\
		                                 & $\|\mathbf{e}-\eP_{h}\|_{L^{2}}$ & 0.56 & 0.72 & 0.84 \\
		                                 & $\|\mP-\mP_{h}\|_{L^{2}}$ & 0.64 & 0.69 & 0.97 \\
		\hline
		\multirow{3}{*}{Fully stabilized} & $\|\mathbf{s}-\sP_{h}\|_{L^{2}}$ & 0.56 & 0.72 & 0.84 \\
		                                 & $\|\mathbf{e}-\eP_{h}\|_{L^{2}}$ & 0.55 & 0.73 & 0.82 \\
		                                 & $\|\mP-\mP_{h}\|_{L^{2}}$ & 0.55 & 0.73 & 0.89 \\
		\hline
		\multirow{3}{*}{Minimally stabilized} & $\|\mathbf{s}-\sP_{h}\|_{L^{2}}$ & 0.55 & 0.73 & 0.82 \\
		                                     & $\|\mathbf{e}-\eP_{h}\|_{L^{2}}$ & 0.55 & 0.72 & 0.82 \\
		                                     & $\|\mP-\mP_{h}\|_{L^{2}}$ & 0.58 & 0.69 & 0.83 \\
		\bottomrule
	\end{tabular}
	\label{tab:vector_errors}
\end{table}


\newpage
\subsection{Numerical convergence under data assignation}

As a final numerical experiment we focus on the mesh convergence under data assignation. To this and, let $\Omega \subset \mathbb{R}^2$ be a bounded Lipschitz domain, discretized by a geometrically conformal and shape-regular finite element mesh $\mathcal{T}_h = \{K_i\}_{i=1}^{N_e}$. The objective is to assess the numerical consistency of the discrete problem with respect to the quantity of available experimental data. For this purpose, we consider a structured partition of the computational domain, into $ N_d $ non-overlapping rectangular subdomains $ \{\Omega_k\}_{k=1}^{N_{sd}} $, such that:

\begin{equation*}
	\overline{\Omega} = \bigcup_{k=1}^{N_d} \Omega_k, \quad \text{with} \quad \text{int}(\Omega_k) \cap \text{int}(\Omega_\ell) = \emptyset \quad \text{for } k \neq \ell.
\end{equation*}

Each subdomain $ \Omega_k $ is centered at $ x_k^c \in \Omega $, and associated with experimental data $ (\etP_k, \stP_k) $. The experimental data, in this case, is based on a manufactured solution we are proposing to work with. Specifically, we admit the smooth field $\uP$ as being: 

\begin{equation*}
	\uP: \Omega \rightarrow \mathbb{R}, \quad \uP = \sin(\pi x)\sin(\pi y),
\end{equation*}

whose gradient is simply given by:

\begin{equation*}
	\eP: \Omega \rightarrow \mathbb{R}^{2}, \quad \eP = \left[\pi \cos(\pi x)\sin(\pi y), \pi \sin(\pi x)\cos(\pi y) \right]^{\top}
\end{equation*}

The exact flux is defined as the negative of the gradient, that is:

\begin{equation*}
	\sP: \Omega \rightarrow \mathbb{R}^{2}, \quad \sP = \left[-\pi \cos(\pi x)\sin(\pi y), -\pi \sin(\pi x)\cos(\pi y)\right]^{\top}
\end{equation*}

Moreover, we set the Lagrange multipliers to zero:

\begin{equation*}
	\begin{aligned}
		&\lP: \Omega \rightarrow \mathbb{R}, \quad  \lP  = 0\\
		&\mP: \Omega \rightarrow \mathbb{R}^{2}, \quad \mP = \left[0,0\right]^{\top}
	\end{aligned}
\end{equation*}

Finally, one might conclude that the source term can be written as:

\begin{equation*}
	q: \Omega \rightarrow \mathbb{R}, \quad q = \left( \zeta + 2\pi^2 \right)\sin(\pi x)\sin(\pi y)
\end{equation*}

Due to the purposes of the experiment, we opted to consider only pure Dirichlet boundary conditions together with $ \text{P}1 $-elements ($ k=0 $) for both natural and equal-order formulations. 

In this setting, let $ \{x_K\} $ denote the centroids of the mesh elements $ K \in \mathcal{T}_h $, and let $ \{x_k^c\}_{k=1}^{N_d} \subset \Omega $ be the centers of the subdomains $ \{\Omega_k\} $, where each $ x_k^c $ is associated with an experimental data pair $ (\etP_k, \stP_k) $. We then define a piecewise-constant approximation of the data over the mesh by assigning to each element the data associated with its nearest sampling center:

\begin{equation*}
	(\etP_h, \stP_h)|_K := (\etP_k, \stP_k) \quad \text{where} \quad k = \arg\min_{j} \| x_K - x_j^c \|.
\end{equation*}

For the equal-order formulation, the assignation is instead performed at the mesh nodes $ \{x_a\} \subset \Omega $, yielding:

\begin{equation*}
	(\etP_h, \stP_h)(x_a) := (\etP_k, \stP_k) \quad \text{where} \quad k = \arg\min_{j} \| x_a - x_j^c \|.
\end{equation*}


\begin{figure}[h!]
	\centering
	\caption{Mesh convergence for the the scalar fields $ \uP_{h} $ and $ \lP_{h} $ in the $ \|\cdot\|_{L^{2}} $ and $ \|\cdot\|_{H^{1}} $ and vector fields $ \eP_{h} $, $ \sP_{h} $ and $ \mP_{h} $ in $ \|\cdot\|_{L^{2}} $ under natural formulation across different number of subdomains $ N_{d} $}
	%\fontsize{11pt}{10pt}\selectfont 
	\def\svgwidth{1.0\columnwidth}
	\input{Figures/PDFs/quadricula_convergence_natural_linear.pdf_tex}
	\label{fig:quadricula_convergence_natural_linear}
\end{figure}


\vfill
\begin{remark}
	Notably, since the dual fields are identically null, the experimental data pair $(\etP, \stP)$ coincides with the exact fields $(\eP, \sP)$ we would deliver to solve the problem. 
\end{remark}

\begin{figure}[h!]
	\centering
	\caption{Unstabilized}
	%\fontsize{11pt}{10pt}\selectfont 
	\def\svgwidth{1.0\columnwidth}
	\input{Figures/PDFs/quadricula_convergence_equal_order_unstabilized_linear.pdf_tex}
	\label{fig:quadricula_convergence_equal_order_unstabilized_linear}
\end{figure}

\begin{figure}[h!]
	\centering
	\caption{Fully stabilized}
	%\fontsize{11pt}{10pt}\selectfont 
	\def\svgwidth{1.0\columnwidth}
	\input{Figures/PDFs/quadricula_convergence_equal_order_fully_stabilized_linear.pdf_tex}
	\label{fig:quadricula_convergence_equal_order_fully_stabilized_linear}
\end{figure}


\begin{figure}[h!]
	\centering
	\caption{Minimally stabilized}
	%\fontsize{11pt}{10pt}\selectfont 
	\def\svgwidth{1.0\columnwidth}
	\input{Figures/PDFs/quadricula_convergence_equal_order_minimally_stabilized_linear.pdf_tex}
	\label{fig:quadricula_convergence_equal_order_minimally_stabilized_linear}
\end{figure}


% ====== DO NOT DELETE THIS ====== %

%\begin{table}[h!]
%	\centering
%	\begin{tabular}{l l l l l}
%		Formulation & Norms & $\gamma = \frac{\pi}{4}$ & $\gamma = \frac{\pi}{2}$ & $\gamma = \frac{3\pi}{4}$ \\
%		\hline
%		\mPltirow{4}{*}{Natural}				 & $\|u-\uP_{h}\|_{L^{2}}$ & 1.135  & 1.438  & 1.646  \\
%		                                 & $\|u-\uP_{h}\|_{H^{1}}$ & 0.5699 & 0.7171  & 0.8139 \\
%		                                 & $\|\lP-\lP_{h}\|_{L^{2}}$ & 1.217  & 1.684 & 1.965  \\
%		                                 & $\|\lP-\lP_{h}\|_{H^{1}}$ & 0.9898  & 0.9941 & 0.9913  \\
%		\hline
%		\mPltirow{4}{*}{Unstabilized}    & $\|u-\uP_{h}\|_{L^{2}}$ & 1.427 & 1.654 & 1.613 \\
%		                                 & $\|u-\uP_{h}\|_{H^{1}}$ & 0.4591 & 0.5946 & 0.6426 \\
%		                                 & $\|\lP-\lP_{h}\|_{L^{2}}$ & 1.161 & 1.301 & 2.161 \\
%		                                 & $\|\lP-\lP_{h}\|_{H^{1}}$ & 0.4593 & 0.296 & 1.243 \\
%		\hline
%		\mPltirow{4}{*}{Fully stabilized} & $\|u-\uP_{h}\|_{L^{2}}$ & 1.204 & 1.635 & 1.843 \\
%		                                 & $\|u-\uP_{h}\|_{H^{1}}$ & 0.5702 & 0.7178 & 0.8148 \\
%		                                 & $\|\lP-\lP_{h}\|_{L^{2}}$ & 1.102 & 1.39 & 1.817 \\
%		                                 & $\|\lP-\lP_{h}\|_{H^{1}}$ & 0.6008 & 0.824 & 0.9891 \\
%		\hline
%		\mPltirow{4}{*}{Minimally stabilized} & $\|u-\uP_{h}\|_{L^{2}}$ & 1.155 & 1.514 & 1.817 \\
%		                                     & $\|u-\uP_{h}\|_{H^{1}}$ & 0.567 & 0.7261 & 0.8138 \\
%		                                     & $\|\lP-\lP_{h}\|_{L^{2}}$ & 1.108 & 1.479 & 1.93 \\
%		                                     & $\|\lP-\lP_{h}\|_{H^{1}}$ & 0.597 & 0.7779 & 0.9834 \\
%		\hline
%	\end{tabular}
%	\label{tab:condkjfs}
%	\caption{Convergence rates for the scalar variables $u$ and $\lP$ under different stabilization schemes.}
%\end{table}

%\begin{table}[h!]
%	\centering
%	\begin{tabular}{l l c c c}
%		Formulation & Norms & $\gamma = \frac{\pi}{4}$ & $\gamma = \frac{\pi}{2}$ & $\gamma = \frac{3\pi}{4}$ \\
%		\hline
%		\mPltirow{3}{*}{Natural}    & $\|\mathbf{s}-\sP_{h}\|_{L^{2}}$ & 0.5676 & 0.7252 & 0.8269  \\
%		                                 & $\|\mathbf{e}-\eP_{h}\|_{L^{2}}$ & 0.5698 & 0.7171 & 0.8138 \\
%		                                 & $\|\mP-\mP_{h}\|_{L^{2}}$ & 0.5817 & 0.7298 & 0.8652 \\
%		\hline
%		\mPltirow{3}{*}{Unstabilized}    & $\|\mathbf{s}-\sP_{h}\|_{L^{2}}$ & 0.5529 & 0.6674 & 0.8582 \\
%		                                 & $\|\mathbf{e}-\eP_{h}\|_{L^{2}}$ & 0.5556 & 0.7233 & 0.8379 \\
%		                                 & $\|\mP-\mP_{h}\|_{L^{2}}$ & 0.6398 & 0.6906 & 0.9706 \\
%		\hline
%		\mPltirow{3}{*}{Fully stabilized} & $\|\mathbf{s}-\sP_{h}\|_{L^{2}}$ & 0.5601 & 0.7196 & 0.8416 \\
%		                                 & $\|\mathbf{e}-\eP_{h}\|_{L^{2}}$ & 0.5495 & 0.7262 & 0.821 \\
%		                                 & $\|\mP-\mP_{h}\|_{L^{2}}$ & 0.5543 & 0.7329 & 0.8948 \\
%		\hline
%		\mPltirow{3}{*}{Minimally stabilized} & $\|\mathbf{s}-\sP_{h}\|_{L^{2}}$ & 0.5496 & 0.7266 & 0.8239 \\
%		                                     & $\|\mathbf{e}-\eP_{h}\|_{L^{2}}$ & 0.5537 & 0.7209 & 0.8224 \\
%		                                     & $\|\mP-\mP_{h}\|_{L^{2}}$ & 0.5820 & 0.6921 & 0.8324 \\
%		\hline
%	\end{tabular}
%	\label{tab:vector_errors}
%	\caption{Convergence rates for the vector variables $\mathbf{s}$, $\mathbf{e}$, and $\mP$ under different stabilization schemes.}
%\end{table}




%=============================================

\section{Conclusions}

This work presents the finite element discretization of
the optimality conditions that appear when solving 
an optimization problem in the emerging field of 
Data-Driven computational mechanics.
We consider the case of the balance law that corresponds to 
the prototypical diffusion-reaction problem and seek for 
a potential field $u$ (e.g., a concentration or a temperature), 
its gradient and the flux $(\eP,\sP)$. In the context of the 
Data-Driven paradigm, the constitutive behavior, relating
gradients and fluxes can be thought as given by functions
$(\etP,\stP)$ originated from experimental data and, as consequence,
do not necessarily satisfy the conservation law. The data-driven
problem, then, consist in finding the fields $(u,\eP, \sP)$ that are
closest to the data, subject to the compatibility and equilibrium
constraints.

We first analyze the well-posedness of such problem in the continuous setting
and second, we propose a family of finite element discretizations to approximate
it. One the one hand, a natural formulation, in which the scalar fields are 
approximated by continuous piecewise polynomial functions and the vector fields
by discontinuous piecewise polynomials with one degree less. On the other hand,
an equal-order formulation with and without stabilizations, that allow to 
use the same continuous interpolation for all fields. 

The convergence properties of the discrete formulations were assessed 
by conducting numerical experiments in 2D and computing the error
with respect to manufactured solutions, in both convex and non-convex
geometries. In all these experiments the data fields $(\etP,\stP)$ were
assumed to be given and available at any point. However, in a final numerical 
experiment we also tested the consistency of the formulation with respect
to the amount of available data points, which is a more representative
case of interest in data-driven applications.

This work can be seen as a first step towards the resolution 
of the full optimization problem that emerges in DDCM, in which
the assignment of the experimental data over the domain 
needs also to be determine. In this sense,
many possible directions can be taken to solve this mixed-integer
nonlinear program (MINLP), among them, the alternating direction schemes, 
in which the data is first distributed based on a current solution, 
and in the sequence, the optimality conditions are solved, so as to update
the solution and repeat the process iteratively. Extensions to 
other physical problems, such as elasticity and porous media
flows is also the subject of ongoing work.
 

%=============================================

% --- References ---
\bibliographystyle{plain}
\bibliography{bibliography}

\end{document}
